\documentclass[11pt,%pdflatex,
dvipdfmx,cjk]{beamer}

%\AtBeginDvi{\special{pdf:tounicode EUC-UCS2}}

\usepackage{amsmath,amssymb,amscd,graphicx}
\usepackage{array}
%\usepackage[pdflatex]{pict2e}
\usepackage{amsthm,color,xcolor}
\usepackage{bm}
\usepackage{framed}
\usepackage{arydshln}

\usepackage{bxdpx-beamer}
\usepackage{pxjahyper}
\usepackage{minijs}
\usepackage{mathtools}
%\def\pdfliteral#1{\special{pdf:content #1}}
%\def\languagename{hoge}

%\usepackage[normal]{fixdif}
%\resetdfont{\mathsf}
\usepackage{comment}

\newcommand{\rr}{\mathbf{R}}
\def\inner<#1>{\langle #1 \rangle}

\makeatletter
\def\pgfutil@insertatbegincurrentpagefrombox#1{%
  \edef\pgf@temp{\the\wd\pgfutil@abb}%
  \global\setbox\pgfutil@abb\hbox{%
    \unhbox\pgfutil@abb%
    \hskip\dimexpr2in-2\hoffset-\pgf@temp\relax% changed
    #1%
    \hskip\dimexpr-2in-2\hoffset\relax% new
  }%
}
\setbeamertemplate{footline}[frame number]

\makeatother

\usetheme{Frankfurt}
\setbeamertemplate{navigation symbols}{}
\title{Kepler problem}
\author{工藤 堅太}
\institute{立命館大学理工学部数理科学科4回}
%\vspace{5mm}Based on a joint project \\
%with Wolfram Bauer \& Abdellah Laaroussi (Leibniz Univ. Hannover)}
\date{\footnotesize{立命館大学理工学部数理科学科幾何学系卒業研究発表会 \\2025年1月27日}}

\begin{document}
\begin{frame}\frametitle{}
\titlepage
\end{frame}


\begin{frame}
\frametitle{内容}
\tableofcontents
\end{frame}

\section{導入}
\begin{frame}
\S 1. 導入
\end{frame}

\begin{frame}
\frametitle{\(n\)次元Kepler問題}
\(n\)次元Kepler問題：
\[
	\ddot{x} = - \frac{x}{\|x\|^{n}}, \qquad x \in \rr^{n} \setminus \{0\}.
\]
\(y = \dot{x}, H(x,y) = \dfrac{\|y\|^{2}}{2} - \dfrac{1}{\|x\|}\)とおくと
\[
	\left\{
	\begin{alignedat}{2}   
	\dot{x} &= y &&= \dfrac{\partial H}{\partial y}(x,y) \\
	\dot{y} &= - \dfrac{x}{\|x\|^{3}} &&= - \dfrac{\partial H}{\partial x}(x,y). 
	\end{alignedat}
	\right.
\]
%このような形の方程式を\textbf{Hamilton関数}\(H\)が定める\textbf{Hamilton方程式}といい，
%\([\partial H/\partial y, \partial H/\partial x]\)を
%\(H\)が定める\textbf{Hamiltonベクトル場}という．
\begin{block}{}
\textbf{目標：}
\(n\)次元Kepler問題の流れと\(\dot{\mathbf{S}}^{n}\)上の測地流が
``同じ''であることを示す．
\end{block}
\end{frame}

\begin{frame}
\frametitle{シンプレクティック多様体}
\begin{block}{定義（シンプレクティック形式）}
$M$：滑らかな多様体．$M$上の2形式$\omega \in \Omega^{2}(M)$が
$M$上のシンプレクティック形式とは，次の2つ条件をみたすこと：
\begin{itemize}
\item 
\(\omega\)は非退化 i.e. 各点\(p \in M\)で\(\omega_{p}: T_{p}M \times T_{p}M \to 
T_{p}^{*}M\)が非退化．

\item
\(\omega\)は閉 i.e. \(\mathsf{d}\omega = 0\)．
\end{itemize}
\end{block}
滑らかな多様体\(M\)と\(M\)上のシンプレクティック形式\(\omega\)の対\((M, \omega)\)を
シンプレクティック多様体という．

\begin{block}{定義（シンプレクティック同相写像）}
$(M_{1}, \omega_{1}), (M_{2},\omega_{2})$：シンプレクティック多様体．微分同相写像
$F\colon M_{1} \to M_{2}$がシンプレクティック同相写像とは，
\[F^{*}\omega_{2} = \omega_{1}.\]
\end{block}
\end{frame}

\begin{frame}
\frametitle{Hamiltonベクトル場}
\begin{block}{命題（Hamiltonベクトル場が一意存在）}
\((M, \omega)\)：シンプレクティック多様体．\(\forall H \in \mathcal{C}^{\infty}(M),
\exists! X \in \mathfrak{X}(M)\) s.t. 
\[
\iota_{X}\omega = - \mathsf{d}H.
\]
%ただし，\(\iota\)は
\end{block}
\(X\)をHamilton関数\(H\)が定めるHamiltonベクトル場といい，\(X_{H}\)と表す．
\end{frame}

\begin{frame}
\frametitle{第1積分}
\begin{block}{定義（第1積分）}
\((M, \omega)\)：シンプレクティック多様体，
\(H\colon M \to \mathbf{R}\)：\(M\)上のHamilton関数, 
\((\varphi_{t}^{X_{H}})_{t \in \mathbf{R}}\)：$X_{H}$の流れ．
滑らかな関数\(G: M \to \rr\)が\((\varphi_{t}^{X_{H}})_{t \in \mathbf{R}}\)の第1積分とは
\[
\frac{\mathsf{d}}{\mathsf{d}t}G(\varphi_{t}^{X_{H}}(x)) = 0, 
\qquad \forall x \in M, t \in \mathbf{R}.
\]
をみたすこと．
\end{block}
\(c \in \mathbf{R}\)が\(G\)の正則値のとき，レベル集合\(G^{-1}(c)\)上に，
力学系\(\varphi_{t}^{X_{H}}|_{G^{-1}(c)}\)が誘導される． %滑らかな場合

\textbf{Remark:}\\
Hamilton関数\(H\)は必ず\((\varphi_{t}^{X_{H}})_{t \in \mathbf{R}}\)の第1積分となる．
\end{frame}

\begin{comment}
\begin{frame}
\frametitle{In local coordinates}
\(M\)：滑らかな多様体．\(M\)の上の余接束\(\pi\colon T^{*}M \to M\)上の自然な1形式
\(\theta \in \Omega^{1}(T^{*}M)\)を\(\xi \in T^{*}M, v \in T_{\xi}(T^{*}M)\)に対して，
\[
\theta_{\xi}(u) = \xi(\mathsf{d}\pi_{\xi}(u)) 
\]
で定める．\(\theta\)を用いて，
\(T^{*}M\)上の自然な2形式\(\omega \in \Omega(T^{*}M)\)を
\(\omega = \mathsf{d}\theta\)
で定めると，\(\omega\)は\(T^{*}M\)上のシンプレクティック形式となる．
\
この\(\omega\)のもとで，Hamilton関数\(H \colon T^{*}M \to \mathbf{R}\)が
定めるHamilton方程式は\(T^{*}M\)の局所座標系\((q,p)\)を用いて，
\[
\begin{cases}
\dot{q} &=  \dfrac{\partial H}{\partial p}(q,p) \\
\dot{p} &=   - \dfrac{\partial H}{\partial q}(q,p)
\end{cases}
\]
と表せる．
\end{frame}
\end{comment}

\begin{frame}
\frametitle{余接束上の自然なシンプレクティック形式}
\(U \subset \rr^{n}\)：開集合．\(T^{*}U = U \times \rr^{n}\)上の
自然なシンプレクティック形式\(\omega \in \Omega^{2}(T^{*}U)\)を
\(U, \rr^{n}\)の標準座標\((q^{1}, \dots, q^{n}), (p_{1}, \dots, p_{n})\)を用いて
次で定める：
\[
	\omega \coloneqq \sum_{i = 1}^{n} \mathsf{d}p^{i} \wedge \mathsf{d}q_{i}.
\]
この\(\omega\)のもとで，Hamilton関数\(H \colon T^{*}U \to \mathbf{R}\)が
定めるHamilton方程式は
\[
\begin{cases}
\dot{q} &=  \dfrac{\partial H}{\partial p}(q,p) \\
\dot{p} &=   - \dfrac{\partial H}{\partial q}(q,p)
\end{cases}
\]
と表せる．
\end{frame}

\section[Geodesics on homogeneous spaces]{Riemannian geodesic flows on homogeneous spaces}
\begin{frame}
\S 2. Riemannian geodesic flows on homogeneous spaces
\end{frame}

\begin{frame}
\frametitle{\(n\)次元Kepler問題}
n次元Kepler問題は\(T^{*}(\rr^{n} \setminus \{0\}) = 
(\rr^{n} \setminus \{0\}) \times \rr^{n}\)上の
Hamilton関数\(H\colon(\rr^{n} \setminus \{0\}) \times \rr^{n} \to \rr\)
\[
	H(q,p) = \frac{1}{2}\|p\|^{2} - \frac{1}{\|q\|}
\]
が定めるHamilton力学系
\[
\begin{cases}
	\dot{q} &=  p \qquad \\
	\dot{p} &= - \frac{q}{\|q\|^{3}} 
\end{cases}
\]
と思える．
\end{frame}

\begin{frame}
\frametitle{\(\mathbf{S}^{n}\)上の測地流}
\(T^{*}\rr^{n+1} = \rr^{n + 1} \times \rr^{n + 1}\)の上の\\
Hamilton関数\(\Phi:\rr^{n + 1} \times \rr^{n + 1} \to \rr\)
\[
	\Phi(\xi, \eta) = \frac{1}{2}(\|\xi\|^{2}\|\eta\|^{2} - \inner<\xi, \eta>)
\]
が定めるHamilton力学系
\[
\begin{cases}
	\dot{\xi} &=  \|\xi\|^{2}\eta -  \inner<\xi, \eta>\xi \\
	\dot{\eta} &= -\|\eta\|^{2}\xi + \inner<\xi, \eta>\eta
\end{cases}		
\]
を考える．この力学系は次の関数を第1積分にもつ：
\[
	\|\xi\|^{2}, \|\eta\|^{2}, \inner<\xi, \eta>.
\]
\end{frame}

\begin{frame}
\frametitle{\(\mathbf{S}^{n}\)上の測地流}
\(\Phi\)が定めるHamilton力学系は，\(\|\xi\|^{2}, \inner<\xi, \eta>\)を第1積分にもつことから
\(T\mathbf{S}^{n} = \{(\xi,\eta) \in \rr^{n + 1} \times \rr^{n + 1} \mid 
\|\xi\| = 1, \inner<\xi, \eta> = 0\}\)上のHamilton力学系
\[
\begin{cases}
	\dot{\xi} &=  \eta \\ 
	\dot{\eta} &= -\|\eta\|^{2}\xi
\end{cases}	
\]
を誘導する．この力学系を\(\mathbf{S}^{n}\)上の\textbf{測地流}という．
\(\|\eta\| = \lambda\)とおくと，この方程式の任意の解は
\[
	\xi(t) = \sin(\lambda t)a + \cos(\lambda t)b, \qquad \exists! a,b \in \rr^{n+1}
\]
と表せる．\(\inner<\xi(t), \xi(t)> = 1, \inner<\xi(t), \dot{\xi}(t)> = 0\)に注意すると，
\(\|a\| =1, \|b\| =1, \inner<a,b> = 0\)であり，
解軌道は\(\mathbf{S}^{n}\)上の大円となることがわかる．
\end{frame}

\begin{frame}
\frametitle{位相同値}
\begin{block}{Definition（位相同値写像）}
	\(M,N\)：滑らかな多様体．
	\((\varphi_{t})_{t \in \rr}, (\psi_{t})_{t \in \rr}\)：\(M, N\)上の力学系．
	同相写像\(h\colon M \to N\)が\((\varphi_{t})_{t \in \rr}\)と\((\psi_{t})_{t \in \rr}\)
	の間の\textbf{位相同値写像}であるとは，
	ある連続写像\(s\colon \rr \times M \to \rr\)が存在して，
	\begin{itemize}
		\item
		各点\(p \in M\)で\(s(\cdot, p)\colon \rr \to \rr\)は単調増加かつ全射．
		
		\item
		\((h \circ \varphi_{s(t,p)})(p) = (\psi_{t} \circ h)(p)\ \forall t \in \rr, p \in M.\)
	\end{itemize}
	をみたすこと．
\end{block}
\((\varphi_{t})_{t \in \rr}\)と\((\psi_{t})_{t \in \rr}\)の間に位相同値写像が存在するとき，
\((\varphi_{t})_{t \in \rr}\)は\((\psi_{t})_{t \in \rr}\)と\textbf{位相同値}であるという．

\textbf{Remark:}\\
位相同値写像を通して第1積分，平衡点，周期軌道などはうつり合う．
\end{frame}

\begin{frame}
\begin{block}{Theorem（Moser）}
負のエネルギー面\(H^{-1}(E)\)\((E < 0)\)に制限した\(n\)次元Kepler問題の流れは
\(\dot{\mathbf{S}}^{n}\)上の測地流に微分同相写像
\(u\colon T_{1}\dot{\mathbf{S}}^{n} \to H^{-1}(E)\)
\[
\begin{cases}
	p_{j} &= \dfrac{\xi}{1 - \xi_{0}}\sqrt{-2E} \\
	-q_{j} &= ((1 - \xi_{0})\eta_{j} + \eta_{0}\xi_{j})(-2E)^{-1}
\end{cases}
\]
と時間変数の滑らかな変更\(\tau: \rr \times H^{-1}(E) \to \rr\)
\[
	\frac{\mathsf{d}t}{\mathsf{d}\tau}(\tau,(q,p)) = \|q\| \cdot (-2E)^{-\frac{3}{2}}
\]
により（微分）位相同値である．
\end{block}
\end{frame}

\begin{frame}
\frametitle{証明の概略}

\end{frame}

\section[証明の詳細]{証明の詳細}
\begin{frame}
\S 3. 証明の詳細
\end{frame}

\begin{frame}
\frametitle{立体射影}
\(L^{+} \coloneqq \{(\xi_{0}, 0, \dots, 0) \mid \xi_{0} \geq 0\}, 
\rr_{+}^{n+1} = \{(x_{0}, \dots ,x_{n}) \in \rr^{n+1} \mid x_{0} > 0\}\)．
``立体射影""
\(a: \rr^{n+1} \setminus L^{+} \to \rr_{+}^{n+1} \)
\[
\begin{cases}
	x_{0} &= \|\xi\| \\
	x_{j} &= \dfrac{\xi_{j}}{\|\xi\| - \xi_{0}}, \qquad j = 1, 2, \dots, n.
\end{cases}
\]
を考える．\(a\)の逆写像は次で与えられる：
\[
\begin{cases}
	\xi_{0} &= x_{0}\dfrac{\|x\|^{2} - 1}{\|x\|^{2} + 1} \\
	\xi_{j} &= x_{0}\dfrac{2x_{j}}{\|x\|^{2} + 1}, \qquad j = 1, 2, \dots, n.
\end{cases}
\]
\end{frame}

\begin{frame}
\frametitle{立体射影}
図
\end{frame}

\begin{frame}
\frametitle{立体射影}
\begin{block}{命題（母関数を用いたシンプレクティック同相写像の構成）}
	\(U, V \subset \rr^{n}\)：開集合．
\end{block}
\end{frame}

\begin{frame}
\frametitle{立体射影}
母関数\(W: \rr_{+}^{n+1} \times \rr^{n+1} \to \rr\)
\[
	W((x_{0,}x), \eta) = x_{0}\dfrac{\|x\|^{2} - 1}{\|x\|^{2} + 1}\eta_{0} + 
			  x_{0}\dfrac{2x_{j}}{\|x\|^{2} + 1}(\sum_{j = 1}^{n}x_{j}\eta_{j})
\]
を用いて，立体射影\(a: \rr^{n+1} \setminus L^{+} \to \rr_{+}^{n+1}\)を
シンプレクティック同相写像
\(u: (\rr^{n+1} \setminus L^{+}) \times \rr^{n+1} \to \rr_{+}^{n+1} \times \rr^{n+1}\)
\[
\begin{cases}
	x_{0} &= \|\xi\| \\
	x_{j} &= \dfrac{\xi_{j}}{\|\xi\| - \xi_{0}} \\
	y_{0} &=  \inner<\xi, \eta>\|\xi\|^{-1} \\
	y_{j} &= (\|\xi\| - \xi_{0})\eta_{j} -(y_{0} - \eta_{0})\xi_{j}, \qquad j = 1, 2, \dots, n.
\end{cases}
\]
にする．
\end{frame}

\begin{frame}
\frametitle{立体射影}
\begin{block}{命題（変換理論）}
	\((M_{1}, \omega_{1}), (M_{2}, \omega_{2})\)：シンプレクティック多様体．
	任意の微分同相写像\(u\colon M_{1} \to M_{2}\)について，
	\(u\)がシンプレクティックであることと,
	\[
		u^{*}X_{H} = X_{u^{*}H}, \qquad \forall H \in \mathcal{C}^{\infty}(M)
	\]
\end{block}
\(F \coloneqq \Phi \circ u^{-1}\)とおくと，
\[
	F((x_{0},x),(y_{0},y)) = \frac{1}{8}(\|x\|^{2} + 1)^{2}\|y\|^{2}
\]
であり，\((\varphi_{t}^{X_{\Phi}})_{t \in \rr}\)と\((\varphi_{t}^{X_{F}})_{t \in \rr}\)は位相同値．

\textbf{Remark:}\\
座標関数\(x_{0}, y_{0}\)は\((\varphi_{t}^{X_{F}})_{t \in \rr}\)の第1積分となっている．
%\[
%	\|\xi\|^{2} = x_{0}^{2}, \qquad \inner<\xi, \eta> = x_{0}y_{0}
%\]
%より，座標関数\(x_{0}, y_{0}\)は\((\varphi_{t}^{X_{F}})_{t \in \rr}\)の第1積分である．
%\(F\)が定めるHamilton方程式は
%\[
%	\dfrac{\mathsf{d}x_{0}}{\mathsf{d}\tau} = 0,\
%	\dfrac{\mathsf{d}x}{\mathsf{d}\tau} = \dfrac{\partial F}{\partial y}(x,y),\
%	\dfrac{\mathsf{d}y_{0}}{\mathsf{d}\tau} = 0,\
%	\dfrac{\mathsf{d}y}{\mathsf{d}\tau} = \dfrac{\partial F}{\partial y}(x,y).
%\]
\end{frame}

\begin{frame}
\frametitle{立体射影}
\(\|\xi\|^{2}, \|\eta\|^{2}, \inner<\xi, \eta>\)は\((\varphi_{t}^{X_{\Phi}})_{t \in \rr}\)の
第1積分であったので
\[
	T_{1}\dot{\mathbf{S}}^{n} = 
 	\{ (\xi, \eta) \in (\rr^{n+1} \setminus L^{+}) \times \rr^{n+1} \mid
	\|\xi\| = 1, \|\eta\| = 1, \inner<\xi, \eta> = 0 \}
\]
上のHamilton関数\(\Phi|_{T_{1}\dot{\mathbf{S}}^{n}}(\xi, \eta) = 1/2\)が定める力学系
\[
\begin{cases}
	\dot{\xi} &=  \eta \\ 
	\dot{\eta} &= -\xi
\end{cases}	
\]
が誘導され，微分同相写像
\(u|_{T_{1}\mathbf{S}^{n}}: T_{1}\mathbf{S}^{n} \to \tilde{F}^{-1}(1/2)\)
によって，Hamilton関数\(\tilde{F}(x,y) = \dfrac{1}{8}(\|x\|^{2} + 1)^{2}\|y\|^{2} \)
が定めるエネルギー面\(\tilde{F}^{-1}(1/2)\)上の力学系と（微分）位相同値である．
\end{frame}

\begin{frame}
\frametitle{Hamilton関数の変更}
\begin{block}{命題（Hamilton関数の変更）}
	\((M, \omega, H)\)：Hamilton力学系，\(c \in \rr\)：\(H\)の正則値．
	滑らかな関数\(\psi\colon \rr \to \rr\)が\(\psi^{\prime}(c) = 1\) を
	満たする．滑らかな関数\(K\colon M \to \rr\)を
	\[
		K \coloneqq \psi \circ H
	\]
	で定めたとき，
	\[
		X_{K} = X_{H}\qquad \text{on}\ H^{-1}(c).
	\]
\end{block}
滑らかな関数\( \psi\colon \rr \to \rr\)を
\(
	\psi(\lambda) \coloneqq \sqrt{2 \lambda}
\)
で定めると，\(\psi\)は\(\psi^{\prime}(1/2) = 1\)を満たし，Hamilton関数\(\tilde{F}\)は
Hamilton関数
\[
	(\psi \circ F)(x,y) = \frac{1}{2}(\|x\|^{2} + 1)\|y\|
\]
と\(F^{-1}(1/2) = (\psi \circ F)^{-1}(1)\)上で\(X_{F} = X_{\psi \circ F}\)となる．
\end{frame}

\begin{frame}
\frametitle{時間変数の変更}
%\begin{block}{命題}
	\((M, \omega, H)\)：Hamilton力学系，\((\varphi_{t})_{t \in \rr}\)を\(X_{H}\)の流れ，
	\(c \in \rr\)：\(H\)の正則値．滑らかな関数\(\lambda\colon M \to \rr_{>0}\)に
	対して，滑らかな写像\(s\colon \rr \times M \to M\)を
	\[
		s(t,x) \coloneqq \int_{0}^{t} \lambda(\varphi_{\tau}(x)) \mathsf{d}\tau \quad
		\text{or} \quad
 		\dfrac{\mathsf{d}s}{\mathsf{d}t}(t, x) = \lambda(x)
	\]
	で定めたとき，各点\(x \in M\)で\(s(\cdot, x)\colon \rr \to \rr\)は
	単調増加かつ全射であるとする．\(M\)上の流れ\((\psi_{t})_{t \in \rr}\)を
	\[
		\psi_{s(t,x)}(x) = \varphi_{t}(x)
	\]
	で定めたとき，\(M\)上の新たなHamilton関数\(L\colon M \to \rr\)を
	\[
		L \coloneqq \lambda^{-1}(H - c)
	\]
	とおくと	
	\[
		X_{L} = \lambda^{-1}X_{H}\qquad \text{on}\ H^{-1}(c).
	\]
%\end{block}
\end{frame}

\begin{frame}
\frametitle{時間変数の変更}
\[
	\frac{\mathsf{d}t}{\mathsf{d}\tau}(\tau,(x,y)) = \|y\|
\]
と時間変数の変更を行なうとHamilton関数\(\psi \circ F\)が定めるHamilton力学系と
Hamilton関数
\[
	\|y\|^{-1}((\psi \circ F)(x,y) -1) = \frac{1}{2}(\|x\|^{2} + 1) - \|y\|^{-1}
\]
が定めるHamilton力学系はエネルギー面\(F^{-1}(1/2)\)上で（微分）位相同値となる． 
\end{frame}

\begin{frame}
最後にシンプレクティック変換\(\rr^{n} \times \rr^{n} \to \rr^{n} \times \rr^{n}\)
\[
	q = -y, \qquad p = x
\]
により\(n\)次元Kepker問題のHamiltonianであるHを得る：
\[
	\dfrac{1}{2}(\|p\|^{2} + 1) - \frac{1}{\|q\|} = H(q,p) + \dfrac{1}{2}
\]
\end{frame}

\begin{frame}
\frametitle{一般の負のエネルギーの場合}

\end{frame}


\appendix
\section*{aa}
\begin{frame}
	\frametitle{今後の展望}
	\begin{itemize}
		\item 
		\(n\)次元Kepler問題の可積分性
	
		\item
		特異点
	
		\item
		\(\mathsf{SO}(n+1)\)対称性
	\end{itemize}
	\end{frame}
	
	%\section{References}
	\begin{frame}
	\frametitle{参考}
	{\footnotesize
	\begin{thebibliography}{999999}
	\beamertemplatetextbibitems
	\bibitem
		[K]{Kawazoe_2024}
		 Kawazoe, 
		 Vector fields on spheres, 
		 Ann. Math., \textbf{75}(3), 603-632, 1962. 
	\end{thebibliography}
	}
	\end{frame}
\end{document}
